\begin{abstract}
SO$(N)$ Chern--Simons 理論の彩色 Kauffman 不変量の large $N$ 双対（Mariño）を、Aganagic--Schaeffer の orientifold の refined 位相的弦の枠組みで refined にする。
SO$(2N)$ の refined Chern--Simons 不変量から、双対 Cauchy 核（t-brane）と、行列式と荷電共役で不変な規格化を用いて、refined のクロスキャップ部分 $g = f - \tfrac12 f^{\cov}$ を取り出す。

主な結果は次の 5 つである。
\begin{enumerate}
\item \textbf{構造}：M 理論の指数で、orientifold の一粒子状態を $\tilde I$ で入れ替わる対・$\tilde I$ で不変な状態・凍結した状態に分けると、全 $d$ の構造 $g = \tfrac12 E/(1+t) + C$ が従う。$E$ と $C$ の分離は指数からは決まらない。閉弦では、SO$(2N)$ の分配関数の積表示から作った自由エネルギーが、$N$ の 1 次式を除いて Aganagic--Schaeffer の式と一致することを確かめた。また、Avetisyan--Mkrtchyan の「Klein 瓶」の寄与が、この分解の並べ替えであることを示す。
\item \textbf{計算}：DAHA-Jones 多項式を一般の $t$ で直接計算し、refined SO$(2N)$ の trefoil の 2 箱・3 箱と $T(2,5)$ の 2 箱を一般の $N$ で求める。大きな $N$ のクロスキャップ部分は整数係数の Laurent 多項式で、$(1+t)$ の極は Adams の予言どおり $p_2$ 方向にだけ現れる。$t = q$ では Mariño の結果を全係数で再現する。
\item \textbf{普遍因子}：refined だけの部分 $\hat g - \hat g|_{t=q}(\lambda \to \lambda (q/t)^{1/2})$ は、[1]・2 箱・3 箱・Hopf 絡み目のすべてで $(\bt^2 - 1)(\ba^2 - \bq^2)$ で割り切れる。Kähler 変数のずれの向きは一意に決まる。
\item \textbf{正値性}：トーラス結び目の [1] では、$G$ 展開の後に符号がそろう判定式が成り立つ。多箱では符号がそろわない。この障害を、凍結した状態の多重項の基底と $E$・$C$ の分離として定式化する。$E$ と $C$ の分離が指数から決まらないことを示し、composite の DAHA の値からも状態の数が得られなかったことを報告する。多箱の $c = 1$ 部分が $\sigma$ 対称でない理由を、unrefined の段階での $c$ の取り違え（hatted では $c = 1$ は $\lambda^\beta$ の $\beta \equiv |R| - 1 \pmod 2$ の部分）と、refined に固有の $\sigma$ 奇部分（$(\bt^2-1)(\ql - \ql^{-1})$ で割り切れる）に分ける。
\item \textbf{Sp$(2N)$}：C$_N$ 型の DAHA を実装し、Avetisyan--Mkrtchyan の refined Sp$(2N)$（$u = t^2$）の結び目不変量を初めて計算した。[1] では $g^{\mathrm{Sp}}(\ba) = -g^{\SO}(-\ba)$（$c = 1$ の符号反転が refined でも成り立つ）。2 箱では彩色の $W_R$ の段階で $W^{\mathrm{Sp}}_R(\ba) = (-1)^{|R|}W^{\SO}_R(-\ba)$ が成り立ち、クロスキャップ部分の差は $p_2$ の定数項と Adams 項に由来する閉じた形 $\pm\psi^2[\delta]$ に尽きる。
\end{enumerate}
\end{abstract}
